\section*{Introduction}

La confidentialité des données numériques repose traditionnellement sur un compromis binaire : la donnée est soit sécurisée (chiffrée et illisible), soit utile (en clair et traitable). Ce paradigme impose que toute manipulation de données sensibles nécessite une phase de déchiffrement, créant ainsi une faille de sécurité critique au moment même où l'information est traitée en mémoire vive. Le \textbf{Chiffrement Totalement Homomorphe} (FHE, pour \textit{Fully Homomorphic Encryption}) a été conçu pour briser cette dualité en permettant l'exécution d'algorithmes arbitraires directement sur des cryptogrammes.

\subsection*{Une quête de trente ans : des prémices à la rupture de Gentry}

L'histoire du FHE débute en 1978, un an seulement après l'invention du RSA. Ronald Rivest, Leonard Adleman et Michael Dertouzos introduisent le concept de « \textit{privacy homomorphisms} ». Ils pressentent déjà l'immense potentiel d'un système où une opération sur les chiffrés correspondrait à une opération sur les clairs. Si certains schémas précoces présentaient des propriétés homomorphes partielles — comme le RSA (homomorphe pour la multiplication) ou le cryptosystème de Paillier en 1999 (homomorphe pour l'addition) — la capacité à combiner \textbf{simultanément} additions et multiplications est restée le « Saint Graal » de la cryptographie pendant trois décennies.

Le verrou technologique saute en 2009 lorsque Craig Gentry, dans sa thèse de doctorat à l'Université de Stanford, publie le premier schéma totalement homomorphe. Sa percée repose sur une idée révolutionnaire : le \textit{bootstrapping}. Puisque chaque opération homomorphe augmente le bruit inhérent au chiffré jusqu'à rendre le déchiffrement impossible, Gentry propose de déchiffrer le texte chiffré \textit{sous sa forme chiffrée} pour réinitialiser ce bruit. Cette preuve d'existence a transformé un rêve théorique en une réalité mathématique, bien que les premiers temps d'exécution se comptaient alors en jours pour une simple opération.

\subsection*{De la théorie à l'efficience : vers une cryptographie multi-facettes}

À la suite des travaux de Gentry, la recherche s'est orientée vers la réduction de la complexité algorithmique et la diversification des structures mathématiques sous-jacentes. On distingue alors plusieurs générations de schémas, passant des réseaux idéaux à des problèmes plus stables comme le \textit{Learning With Errors} (LWE). 

Ce mémoire se propose d'analyser cette évolution à travers trois protocoles emblématiques qui illustrent la maturité actuelle du domaine :

\begin{itemize}
    \item \textbf{Le schéma DGHV (2010) :} Proposé par van Dijk, Gentry, Halevi et Vaikuntanathan, il simplifie la construction de Gentry en remplaçant les réseaux complexes par l'arithmétique des entiers. L'étude de DGHV est pédagogiquement essentielle pour comprendre la gestion du bruit via le problème du diviseur commun approximatif (AGCD).
    
    \item \textbf{Le schéma TFHE (2016) :} Représentant de la troisième génération, le \textit{Torus FHE} se distingue par son optimisation radicale du \textit{bootstrapping}. En travaillant sur le tore réel, il permet d'évaluer des portes logiques (NAND, XOR) en quelques millisecondes, ouvrant la voie à l'exécution de circuits booléens complexes et de microprocesseurs chiffrés.
    
    \item \textbf{Le schéma CKKS (2017) :} Conçu par Cheon, Kim, Kim et Song, ce schéma marque un tournant pragmatique. Contrairement aux schémas exacts, CKKS traite le bruit comme une erreur d'arrondi inhérente au calcul en virgule flottante. Cette approche « approximative » le rend incomparablement plus efficace pour le \textit{Machine Learning} et le traitement de données statistiques à grande échelle.
\end{itemize}

L'enjeu de ce travail sera de démontrer comment ces constructions mathématiques, bien que radicalement différentes dans leurs approches, convergent vers un objectif commun : rendre le traitement de données privées aussi fluide et universel que le calcul en clair.